\section*{第 \ref{ch:g11:lenses-and-eye} 章——透镜、成像与眼睛}

\begin{solution}{exo:g11:lenses-and-eye:1}
$C = 1/0.50 = +2.0$~屈光度；$1/(-0.20) = -5.0$~屈光度；$1/0.0040 = +250$~屈光度。
一块 $+8.0$~屈光度的透镜：$f' = 1/8.0 = \qty{0.125}{m} = \qty{12.5}{cm}$。
\end{solution}

\begin{solution}{exo:g11:lenses-and-eye:2}
凹透镜（$f' < 0$）。不能。对任意物距 $d$，
$\gamma = f'/(f' - d)$（见\cref{exo:g11:lenses-and-eye:15}）
介于 $0$ 与 $1$ 之间，像总是
虚像、正立且\emph{缩小}。印刷字样看起来正立
而又缩小，从不放大。
\end{solution}

\begin{solution}{exo:g11:lenses-and-eye:3}
物体位于 $2f'$ 处；过光心的光线与平行光线相交于
透镜后 \qty{20}{cm} 处。校验：
$1/\overline{OA'} = 1/10 - 1/20 = 1/20$，故
$\overline{OA'} = +\qty{20}{cm}$，$\gamma = 20/(-20) = -1$：实像，
倒立、等大——与图中一致。
\end{solution}

\begin{solution}{exo:g11:lenses-and-eye:4}
$1/\overline{OA'} = 1/8.0 - 1/24 = 1/12$：
$\overline{OA'} = +\qty{12}{cm}$，$\gamma = 12/(-24) = -0.50$。实
像位于透镜后 \qty{12}{cm} 处，倒立，大小为原物的一半。
\end{solution}

\begin{solution}{exo:g11:lenses-and-eye:5}
角膜与晶状体共同构成凸透镜；视网膜
即屏幕。调节改变的是晶状体的焦距（光焦度），
而晶状体与视网膜间的距离不能改变。近
点：在完全调节时清晰可见的最近点——
对标准眼而言为 \qty{25}{cm}。
\end{solution}

\begin{solution}{exo:g11:lenses-and-eye:6}
风景照：感光元件位于焦平面，距透镜 \qty{50}{mm}。在
\qty{3.0}{m} 处：$1/\overline{OA'} = 1/50 - 1/3000 = 59/3000$，故
$\overline{OA'} \approx \qty{50.8}{mm}$。行程约为 \qty{0.8}{mm}。
\end{solution}

\begin{solution}{exo:g11:lenses-and-eye:7}
$1/\overline{OA'} = 1/6.0 - 1/5.0 = -1/30$：
$\overline{OA'} = \qty{-30}{cm}$——虚像、正立，
$\gamma = (-30)/(-5.0) = +6.0$：像为邮票的六倍大，
在透镜同侧 \qty{30}{cm} 处。
\end{solution}

\begin{solution}{exo:g11:lenses-and-eye:8}
$1/f' = 1/40 + 1/40 = 1/20$：$f' = \qty{20}{cm}$；
$\gamma = 40/(-40) = -1$。对称性要求
$\abs{\overline{OA'}} = \abs{\overline{OA}}$，因此
$\abs{\gamma} = 1$；像为实像且倒立，故 $\gamma = -1$：
恰好等大。
\end{solution}

\begin{solution}{exo:g11:lenses-and-eye:9}
$1/\overline{OA'} = -1/25 - 1/50 = -3/50$：
$\overline{OA'} \approx \qty{-16.7}{cm}$，
$\gamma = (-16.7)/(-50) = +0.33$：虚像、正立，像的大小为
原来的三分之一。近视者的眼镜是凹透镜：通过它们看到的一切
——包括从外面看到的佩戴者的眼睛——都是缩小的。
\end{solution}

\begin{solution}{exo:g11:lenses-and-eye:10}
无穷远物体的像在远点处：$f' = \qty{-0.40}{m}$，故
$C = -2.5$~屈光度。山的像位于
透镜前 \qty{40}{cm} 处，恰为远点，放松的眼睛
无需调节即可清晰看到。
\end{solution}

\begin{solution}{exo:g11:lenses-and-eye:11}
$C = 1/\overline{OA'} - 1/\overline{OA} = 1/(-0.80) - 1/(-0.25)
= -1.25 + 4.00 = +2.75$~屈光度。
\end{solution}

\begin{solution}{exo:g11:lenses-and-eye:12}
\emph{1.} $\gamma = \overline{OA'}/\overline{OA} = -50$，得
$\overline{OA'} = -50\,\overline{OA}$；结合
$\overline{OA'} - \overline{OA} = \qty{5.1}{m}$，
$-51\,\overline{OA} = \qty{5.1}{m}$：$\overline{OA} = \qty{-0.10}{m}$，
$\overline{OA'} = +\qty{5.0}{m}$。

\emph{2.} $1/f' = 1/5.0 + 1/0.10 = 10.2$：$f' \approx \qty{9.8}{cm}$。

\emph{3.} $50 \times \qty{24}{mm} = \qty{1.2}{m}$，倒立——这
正是幻灯片要倒着装入的原因。
\end{solution}

\begin{solution}{exo:g11:lenses-and-eye:13}
\emph{1.} $C_\infty = 1/0.017 \approx 59$~屈光度；
$C_{25} = 1/0.017 + 1/0.25 \approx 63$~屈光度。振幅：
$1/0.25 = 4.0$~屈光度。

\emph{2.} 振幅 $= 1/1.0 = 1.0$~屈光度：仅余
四分之一（$25\%$）的标准眼 $4.0$~屈光度。
\end{solution}

\begin{solution}{exo:g11:lenses-and-eye:14}
\emph{1.} $f' = 0.25/3 \approx \qty{0.083}{m} = \qty{8.3}{cm}$。

\emph{2.} $\overline{OA'} = \qty{-25}{cm}$：
$1/\overline{OA} = -1/25 - 3/25 = -4/25$，故
$\overline{OA} = \qty{-6.25}{cm}$，且
$\gamma = (-25)/(-6.25) = +4.0$——优于标称的
$\times 3$（该值假定像位于无穷远）；将像移至
近点时可再增益一个单位。
\end{solution}

\begin{solution}{exo:g11:lenses-and-eye:15}
$1/\overline{OA'} = 1/f' - 1/d = (d - f')/(f'd)$，故
$\overline{OA'} = f'd/(d - f')$，且
$\gamma = \overline{OA'}/(-d) = f'/(f' - d)$。若 $d > f'$：
$\overline{OA'} > 0$（实像），$\gamma < 0$（倒立）。若
$0 < d < f'$：$\overline{OA'} < 0$（虚像），且
$\gamma = f'/(f' - d) > 1$：正立且放大——即放大镜的
用法。
\end{solution}

\begin{solution}{pb:g11:lenses-and-eye:1}
\textbf{1.} $1/\overline{OA'} = 1/20 - 1/60 = 1/30$：
$\overline{OA'} = +\qty{30}{cm}$，$\gamma = -0.50$——实像，倒立，
缩小一半。

\textbf{2.} $\overline{OA'} = +\qty{60}{cm}$，$\gamma = -2$：两个
位置\qty{30}{cm}与\qty{60}{cm}互换了角色，
放大率互为倒数——逆向光线沿同一
路径传播。

\textbf{3.} $1/\overline{OA'} = 1/20 - 1/10 = -1/20$：
$\overline{OA'} = \qty{-20}{cm}$，$\gamma = +2$——虚像，正立，
放大一倍。

\textbf{4.} 物体自无穷远处移来时，像起始于焦平面上；当物体
靠近 $F$ 时，倒立的实像远离透镜并趋于无穷远，
且不断增大；位于 $f'$ 以内时，像变为虚像，正立、放大，位于
透镜同侧。

\textbf{5.} 取 $1/\overline{OA} \approx 0$，则
$1/\overline{OA'} = 1/f'$：$\overline{OA'} = f'$，即焦平面。

\textbf{6.} 位于焦平面上，在透镜后方\qty{50}{mm}处。

\textbf{7.} $1/\overline{OA'} = 1/50 - 1/1000 = 19/1000$：
$\overline{OA'} \approx \qty{52.6}{mm}$；行程 $\approx \qty{2.6}{mm}$。

\textbf{8.} $1/\overline{OA} = 1/55 - 1/50 = -1/550$：最近对焦
距离为\qty{550}{mm} = \qty{55}{cm}。

\textbf{9.} $\gamma = 55/(-550) = -0.10$：面部的像位于
\qty{20}{mm}处，恰好落在\qty{24}{mm}的感光元件之内。

\textbf{10.} $\gamma = -1$ 要求 $\overline{OA'} = -\overline{OA}$，
于是 $2/\overline{OA'} = 1/f'$：物体位于前方\qty{100}{mm}（$= 2f'$）处，
感光元件位于后方\qty{100}{mm}处——这是风景拍摄所需行程的两倍，远超
\qty{55}{mm}的机构行程；微距镜头提供了额外的行程。

\textbf{11.} 比她的近点（\qty{1.0}{m}）更近的物体都会模糊：
伸直手臂可将书页尽量移至距眼\qty{1.0}{m}处——
虽小，却清晰。

\textbf{12.} $C = 1/(-1.0) - 1/(-0.25) = -1.0 + 4.0 = +3.0$~D。

\textbf{13.} $\gamma = (-1.0)/(-0.25) = +4.0$：像大四倍，但
距离也远四倍——视角大小不变；所得收益不在于大小，而在于
\emph{清晰度}，因为像此时恰好位于她的近点处。

\textbf{14.} 要清晰成像，像须位于\qty{1.0}{m}之外：书页由
\qty{25}{cm}处（像位于\qty{1.0}{m}处）移至
$f' = 1/3.0 \approx \qty{0.33}{m}$处（像位于无穷远）。远处
物体经眼镜在其\emph{后方}\qty{33}{cm}处成实像——眼睛无法利用
这样的像：远处景物变得模糊。

\textbf{15.} 因此需要双光镜：镜片下半部分提供阅读所需的屈光度，
上半部分不加矫正（或仅作远距矫正）。

\textbf{16.} $1/\overline{OA'} = 1/5.0 - 1/4.0 = -1/20$：
$\overline{OA'} = \qty{-20}{cm}$，虚像，正立，$\gamma = +5.0$。

\textbf{17.} $1/\overline{OA} = -1/25 - 1/5.0 = -6/25$：
$\overline{OA} \approx \qty{-4.2}{cm}$；$\gamma = (-25)/(-25/6) = +6.0$。

\textbf{18.} 斯通位于 $F$ 处：像位于无穷远，观察时无需调节——
可以连续数小时检视而不致眼睛疲劳。

\textbf{19.} $C_\infty = 1/0.017 \approx 59$~D；
$C_{25} = 1/0.017 + 1/0.25 \approx 63$~D：约 $4$~D 的
调节量。

\textbf{20.} $f' = 1/58.8 \approx \qty{17.0}{mm}$，而
$1/62.8 \approx \qty{15.9}{mm}$，相差约\qty{1.1}{mm}，约 $6\%$——
这正是相机借助\qty{2.6}{mm}的行程所模拟的内置变焦，
而祖母则以 $+3$~D 的镜片加以弥补。
\end{solution}
